\documentclass[10pt,a4paper]{amsart}
\usepackage[margin=19mm]{geometry}
\usepackage{amsmath,amssymb,mathtools}
\usepackage[scr=rsfs,cal=euler]{mathalfa}
\usepackage{xfrac}
\usepackage[unicode,hidelinks,bookmarks=false]{hyperref}
\newcommand{\ie}{i.e.\ }
\newcommand{\cf}{cf.\ }
\newcommand{\resp}{respectively\ }
\newcommand{\Subsection}{\subsection}
\providecommand{\nobreakdash}{\nobreak\hbox{-}}
\setlength{\emergencystretch}{2em}
\multlinegap0pt
\usepackage{amssymb}
\newtheorem{theorem}{Theorem}
\newtheorem{corollary}{Corollary}
\newcommand{\F}{\mathbb{F}}
\title{Algebraic Geometry Codes Approach the Half-Singleton Bound with Constant Field Size}
\author{Neehar Verma, Camilla Hollanti, Razane Tajeddine}
\date{Source: arXiv:2609.05017v1 (CC BY 4.0)}
\begin{document}
\maketitle

\noindent\emph{Source and licence.} Algebraic Geometry Codes Approach the Half-Singleton Bound with Constant Field Size. Version arXiv:2609.05017v1 (CC BY 4.0), distributed by arXiv under the Creative Commons Attribution 4.0 International licence, http://creativecommons.org/licenses/by/4.0/, as recorded in the arXiv licence metadata for this exact version and shown on the abstract page https://arxiv.org/abs/2609.05017v1. Full licence notice: \texttt{licenses/arxiv-2609.05017-CC-BY-4.0.txt}.\par\smallskip

\noindent\textit{Editorial scope.} Counted unit: the article's corollary corollary (source0.tex lines 455--461). The article's complete original proof of it is delivered with it (source0.tex lines 462--489, one \texttt{proof} environment of 28 lines). No mathematics is written, repaired or re-derived: the statement and the proof are the article's own lines, in the article's own notation, and no retained line is substituted or re-flowed.  Retained uncounted as the source-local context the proof needs: lines 343--347.  Nothing was omitted from any retained span, so the package carries no omission record.  No retained line contains a citation, so the wrapper carries no bibliography and no bibliography entry.  No retained line contains a figure environment, an image-inclusion command or drawing code, so no image is delivered and the package declares no asset dependency.  Editorial formatting only, in addition: an AMS \texttt{amsart} preamble that restates the article's own declarations and macros used by the retained lines, namely the environments \texttt{theorem}, \texttt{corollary} on separate counters, together with the standard packages those lines need.  The retained environments print in this document as Theorem 1, Corollary 1 in order of appearance, whereas the article prints the same items with the numbering of its own full document.  The complete native source of the article as distributed in the arXiv e-print for this exact version is bundled unchanged as \texttt{sources/209406/source0.tex}.
\begin{theorem} \label{ultimate}
    Let $C_{\mathcal{F}}(\mathcal{P}) \subseteq \F_q^n$ be an $[n,k,d]$ evaluation code with associated vector space $\mathcal{F} = \langle f_0, f_1, \dotsm , f_{k-1}\rangle_{\F_q}$ of functions (with $f_0 = 1$) and evaluation vector $\mathcal{P} = (P_1, \dots, P_n)$. Let $\ell \geq 2k-1$, then
    $$\Pr[C_{\mathcal{F}}(\mathcal{P}) \textit{ cannot correct $n-\ell$ insdel errors}] \leq \binom{n}{l}^2 \binom{\ell - (n - d + 1)}{\tau} \left(\frac{Z}{|\mathcal{X}| - n + 1}\right)^\tau,$$
    where $\tau  := \ell - k - (n - d) + 1$.
\end{theorem}

\begin{corollary}
    Let $\varepsilon \in (0,1)$ and $R \in (0, 1/2)$. Let $p$ be a prime power such that 
    $$p \geq 1 + \max\left\lbrace \frac{4}{R}, \frac{8}{\varepsilon}\right\rbrace \left(1 +2^{8/\varepsilon}\right),$$
    for every sufficiently large $n$ for which there exists an integer $m$ such that $$ \min\left\lbrace \frac{Rn}{4}, \frac{\varepsilon n}{8}  \right\rbrace \leq p^m \leq \min\left\lbrace \frac{Rn}{2}, \frac{\varepsilon n}{4}\right\rbrace,$$
    and set $q = p^2$. 
    Suppose $\mathcal{P} = (P_1, \dots, P_n)$ is chosen uniformly at random from the set of all $n$-tuples of distinct elements in $\mathcal{P}_q(\mathcal{Y}_m) \setminus P_\infty$, where $\mathcal{Y}_m$ is the $m^{th}$ curve in the GS tower with genus $g_m$. Let $G = (k - 1 + g_m)P_{\infty}$ be a divisor. Then, with probability at least $1-2^{-n}$ the $[n, k = Rn]$ AG code $\mathcal{C}_{\mathcal{L}(G)}(\mathcal{P})$ can correct at least $(1-\varepsilon)n - 2k + 1$ insdel errors.
\end{corollary}

\begin{proof}
    We apply the general Theorem \ref{ultimate} with $\mathcal{X} = \mathcal{P}_q(\mathcal{Y}_m) \setminus \{P_\infty\}$ and the function space as the Riemann--Roch space $\mathcal{F} =\mathcal{L}(G)$.
    For the genus $g_m$ of $\mathcal{Y}_m$ we have the upper bound $g_m \leq 2p^m$ which along with $p^m \leq \frac{Rn}{2}$ implies $$g_m \leq Rn = k. $$
    This forces $\deg(G) = k + g_m - 1 \geq 2g_m - 1$ which consequently tells us by the Riemann--Roch theorem that the function space $\mathcal{L}(G)$ has dimension $\ell(G) = k$. 
    For any function $f \in \mathcal{F}$ we have that $\deg(G)$ bounds the number of zeros, and so we set our zero bound $Z = \deg(G)$. 
    Take $\ell = 2k - 1 + \lfloor \varepsilon n \rfloor$. The assumed inequality $p^m \leq \frac{\varepsilon n}{4}$ tells us that $g_m \leq \frac{\varepsilon n}{2}$ which forces 
    \begin{align*}
        \tau &:= \ell - k - (n-d) + 1 \\
        & \geq 2k - 1 + \lfloor \varepsilon n \rfloor - k - (k-1+g_m) + 1 \\
        & \geq  \lfloor \varepsilon n \rfloor + 1 - g_m \\
        & \geq \left\lceil\frac{\varepsilon n}{2} \right\rceil .
    \end{align*}
    Then, by Theorem \ref{ultimate} we have 
    $$\Pr[C_{\mathcal{F}}(\mathcal{P}) \textit{ cannot correct $n-2k + 1 - \lfloor \varepsilon n \rfloor$ insdel errors}] \leq 2^{3n}\left(\frac{\deg(G)}{|\mathcal{P}_q(\mathcal{Y}_m) \setminus \{P_\infty\}| - n + 1}\right)^{\frac{\varepsilon n}{2}}.$$
    Note that
    \begin{align*}
        \frac{\deg(G)}{|\mathcal{P}_q(\mathcal{Y}_m) \setminus \{P_\infty\}| - n + 1} &= \frac{Rn + g_m - 1}{p^m(p-1) - n + 1} \\
        & \leq \frac{Rn + \frac{\varepsilon n}{2}  - 1}{p^m(p-1) - n + 1} \qquad \text{(by $g_m \leq \frac{\varepsilon n}{2}$)}\\
        & \leq \frac{n}{p^m(p-1) - n + 1} \qquad \text{(by $\varepsilon \in (0,1)$ and $R \in (0, 1/2)$)}\\
        & \leq \frac{p^m \max\left\lbrace \frac{4}{R}, \frac{8}{\varepsilon}\right\rbrace}{p^m(p-1)- p^m\max\left\lbrace \frac{4}{R}, \frac{8}{\varepsilon}\right\rbrace + 1} \qquad \text{(by $p^m \geq \min\left\lbrace \frac{Rn}{4}, \frac{\varepsilon n}{8}\right\rbrace$)} \\
        & \leq \frac{1}{\min\left\lbrace \frac{R}{4}, \frac{\varepsilon}{8}\right\rbrace(p-1)-1}. 
    \end{align*}
    This simplifies the failure probability to 
    $$\Pr[C_{\mathcal{F}}(\mathcal{P}) \textit{ cannot correct $n-2k + 1 - \lfloor \varepsilon n \rfloor$ insdel errors}] \leq 2^{3n}\left(\frac{1}{\min\left\lbrace \frac{R}{4}, \frac{\varepsilon}{8}\right\rbrace(p-1)-1}\right)^{\frac{\varepsilon n}{2}}.$$
    With the consideration $p \geq 1 + \max\left\lbrace \frac{4}{R}, \frac{8}{\varepsilon}\right\rbrace \left(1 +2^{8/\varepsilon}\right)$ this gives us the desired
    $$\Pr[C_{\mathcal{F}}(\mathcal{P}) \textit{ cannot correct $n-2k + 1 - \lfloor \varepsilon n \rfloor$ insdel errors}] \leq 2^{-n}.$$
    Therefore, with probability at least $1-2^{-n}$, the random choice of evaluation points $\mathcal{P} \in \mathcal{X}$ yields a code capable of correcting at least $n-\ell = n-2k+1-\lfloor\varepsilon n\rfloor \geq (1 - \varepsilon)n - 2k + 1$ insdel errors over a field of size $q  = p^2= 2^{O(1/\varepsilon + \log(1/R))}$.
\end{proof}

\end{document}
